\section{结论}

本文研究已知词级对齐条件下非侵入式语言解码的有限等待与滚动预测修正。
通过比较局部词窗与完整上下文模型，我们发现两个自然连续聆听数据集上的局部上下文能够在固定候选词级检索中取得良好表现。
在此基础上，共享滚动模型允许目标词在自身神经窗可用后先行预测，
再利用最多两个后续词窗更新近期预测，从而形成初判、回修与最终固定的连续输出过程。

在现有序列解码协议下，法语 MEG 的滚动修正改善了最终文本质量，中文 EEG 的收益主要体现于 WER 和词位准确率；
有限回修同时带来了可量化的输入等待与输出改写代价。
本文的结论限于已知词级对齐的离线重放设置。
未来可考虑将有限等待和历史预测修正推广到无精确事件对齐的 CTC 类连续神经解码框架；
在连续输入下限制未来帧可见范围，进一步评估预测质量、等待和输出稳定性。
